\section{Formal control model}
\label{sec:formalism}

\subsection{Work, content, and read identity}

Let $R$ be the raw identifiers emitted by retrieval routes. Canonicalization
maps each identifier to a scheme-qualified work key and a rendition version.
Alias-connected keys form a work. Each retrieved rendition also carries a
content digest $c$. Deduplication and overlap are computed over content digests,
not locators, because reminted identifiers must not make two captures appear
independent.

A read receipt is a tuple $(w,c,s,f)$ containing work, content digest,
extraction-schema version, and scientific frame. A candidate is already read
only when all four coordinates match an existing receipt. A changed content
digest, schema, or frame creates a legitimate reread. This rule separates
duplicate processing from processing that is newly required by the scientific
question.

\subsection{Earned route independence}

\begin{definition}[Route capture]
A route capture is $\rho=(k,b,q,C)$, where $k$ is route kind, $b$ is backend
identity, $q$ is the query-derivation process, and $C$ is the captured set of
content identities.
\end{definition}

For two route captures, the controller accepts an independence claim only if
route kind, backend identity, and query derivation differ:
\[
 \operatorname{ind}(\rho_i,\rho_j)
 \Longleftrightarrow k_i\ne k_j\wedge b_i\ne b_j\wedge q_i\ne q_j.
\]
This predicate is a conservative design requirement, not a proof of statistical
independence. Its purpose is to reject obvious shared-generation causes before
overlap is interpreted. Zero observed overlap alone does not establish
independence.

For an accepted pair with capture sizes $n_i,n_j$ and overlap $m_{ij}$, a
two-source diagnostic may report $\widehat N=n_i n_j/m_{ij}$. The diagnostic is
undefined when $m_{ij}=0$. In every case it remains advisory because adaptive
search violates the interpretation of independent capture occasions unless a
separate design establishes it.

\subsection{Route control and task closure}

Each route records attempts, retrieved and novel content identities, resource
use, failure signatures, and backend exhaustion. A priority-ordered policy may
stop on budget, stop on reported exhaustion, suspend on transport failure,
switch after a flat streak, reformulate, or continue. Backend exhaustion is a
claim about the index. Transport failure is a claim about the channel. The
latter suspends the route and preserves an open obligation.

\begin{proposition}[Route-stop non-certification]
No route decision certifies task completeness. Task closure is a separate
decision over the ensemble of declared material obligations.
\label{prop:failclosed}
\end{proposition}

The proposition is an implementation invariant rather than an empirical law.
Its consequence is testable: if a route-stop signal is promoted directly to
task closure, the evaluator should observe premature closure whenever another
route still contains reachable relevant material. Task closure is permitted
only when all declared material routes are discharged or explicitly resolved
by the task contract, no transport failure leaves a material obligation open,
and no advisory coverage diagnostic supplies the authority.

\subsection{Observable error families}

The model defines four evaluator-visible failures. A route-stop false positive
occurs when a stopped route can still reach relevant material. A route-stop
false negative occurs when search continues beyond route exhaustion. Premature
task closure occurs when relevant material remains reachable within the task
contract. Duplicate processing occurs when all four read coordinates repeat;
changed content, schema, or frame instead defines a legitimate reread. Complete
gold is required wherever these definitions depend on the full relevant set.
